\textbf{Who it is for.} Wealth managers and private banks that advise \emph{cautious to balanced private clients},
typically aged 45--65, a few years before or at the start of retirement, with an investment horizon of about 7--10 years
over which they will start drawing on their savings. For these clients a deep loss is the main danger: once withdrawals
begin, assets are sold at low prices and the loss cannot be recovered by waiting (\emph{sequence-of-returns risk}). They
accept a lower return for a smoother path, still want growth above cash and inflation, and want a rule that their adviser
can explain on one page. The strategy is run in US dollars; euro-based clients use the UCITS funds and the currency hedge
of Section~7.

\textbf{Why not simply 60/40?} A 60/40 portfolio is the natural default for such clients, which is why it is our benchmark.
Its weakness is that equities dominate its risk: it lost \res{b6040_ret_2008} in 2008 and \res{b6040_ret_2022} in 2022,
exactly the kind of loss that makes a client sell at the bottom. The strategy is designed for clients who would abandon a
60/40 portfolio in such a crash.

\textbf{Why these six assets.} Each holding protects the client in a different economic scenario: US and international
equities and real estate grow when the economy expands, Treasuries gain in recessions, and commodities and gold hold their
value when inflation surprises. Nobody knows which scenario comes next, so the portfolio holds all of them, lets none of
them dominate its risk, and moves into T-bills when an asset class is in a sustained decline.

\textbf{Who should not buy it.} Clients who measure success against 60/40 or the stock market: the strategy beat 60/40 in
only \res{years_beating_6040} of \res{n_years_cal} calendar years and lags in long bull markets. Investors seeking maximum
long-term growth with a long horizon should hold more equities instead.
